\documentclass[11pt]{article}
\usepackage[T1]{fontenc}\usepackage{lmodern,microtype,amsmath,amssymb,amsthm,booktabs}
\usepackage[margin=1.04in]{geometry}
\usepackage{xcolor}\usepackage[colorlinks=true,allcolors=blue!45!black]{hyperref}
\newtheorem{theorem}{Theorem}\newtheorem{lemma}[theorem]{Lemma}
\newtheorem{corollary}[theorem]{Corollary}\newtheorem{proposition}[theorem]{Proposition}
\newcommand{\Q}{\mathbb Q}\newcommand{\Z}{\mathbb Z}
\title{\bfseries Which exponent digits determine\\ the roots of a sparse $p$-adic polynomial?}
\author{Research note}\date{November 27, 2026}
\begin{document}\maketitle
\begin{abstract}
We give a finite congruence criterion for the entire deep-disk root-count profile of a univariate $p$-adic polynomial space with prescribed exponents, for any sparsity. It produces an exact algorithm even when arbitrarily many exponents have the same residue. We also determine the sharp maximum in the open exponential disk for every $p+2\le t\le 2p+1$: the answer is $p^2+\lfloor(t-p-2)/2\rfloor$. This continues the first jump found in the ProveIt article and answers its question on root counts beyond $p$ terms in an explicit range. The results count geometric roots with multiplicity and the maxima have simple-root witnesses.
\end{abstract}

\section{Disks, supports, and a finite invariant}
Normalize $v=v_p$ by $v(p)=1$, set $v(0)=+\infty$ and take coefficients in any finite extension $K$ of $\Q_p$. Count roots in an algebraic closure with multiplicity. With $r_0=1/(p-1)$ we study
\[
   B_r=\{x:v(x-1)\ge r\}\quad(r>r_0),\qquad B^\circ=\{x:v(x-1)>r_0\}.
\]
A nonzero center $a$ in place of $1$ is covered by $X\mapsto aX$ and relative depth $v(x/a-1)$. A support set $E$ of size $t$ consists of distinct nonnegative integers. Its polynomial space is $V_E(K)=\{\sum_{n\in E}c_nX^n:c_n\in K\}$. Zero coefficients are allowed and the zero polynomial is excluded when counting roots. Let $N_p^\circ(t)$ denote the greatest number of roots in $B^\circ$ for at most $t$ terms, allowing all supports and all coefficient fields. Write $S_p(i)$ for the base-$p$ digit sum and set, for $t\ge2$,
\[
 D=\max_{0\le i<t}S_p(i),\qquad M=\left\lceil\frac{D-1}{p-1}\right\rceil,\qquad q=p^M .
\]
For $t=1$ there are no roots in any of these disks. Our first theorem matches actual polynomials on two supports and works at every depth at once.

\begin{theorem}[Finite congruence for the entire profile]\label{transfer}
Let $E,F$ have size $t$ and $f_E(T)=\prod_{n\in E}(T-n)$. If $f_E\equiv f_F\bmod p^M$ coefficientwise, there is a $K$-linear isomorphism $V_E(K)\to V_F(K)$ under which corresponding polynomials have the same number of roots in every $B_r$ and in $B^\circ$. In particular, the exponent multiset modulo $q$ suffices. The linear isomorphism works for every finite coefficient extension. It is specified over $\Q$ at the center $1$.
\end{theorem}

Multiplicity matters when speaking of an exponent multiset modulo $q$. A stronger assertion than equality of maxima is that the depths $v(x-1)>r_0$ agree as multisets, including any roots at $1$, since the number in every nested disk agrees. The argument gives no correspondence between individual roots outside these disks.

\begin{lemma}[Normalized moments]\label{moment}
For $P=\sum c_nX^n$, put $M_i(P)=\sum_{n\in E} c_n n^i$, with $n^0=1$. Divide the moments by a nonzero initial moment of smallest valuation among $0\le i<t$, and denote the normalized sequence by $m_i$. Every $m_i$ is integral. If $r>r_0$ is rational, the number of roots in $B_r$ is the largest minimizing index of
\[
 v(m_k)+w_r(k),\qquad
 w_r(k)=rk-v(k!)=(r-r_0)k+\frac{S_p(k)}{p-1}.
\]
Every minimizing $k\ge t$ obeys
\begin{equation}\label{digit}
 S_p(k)<D,\qquad v(m_k)<\frac{D-S_p(k)}{p-1}\le M .
\end{equation}
\end{lemma}
\begin{proof}
The initial moment transformation has nonzero Vandermonde determinant. The moments obey the monic recurrence with characteristic polynomial $f_E$, by evaluation at each $n$. Integrality follows from that recurrence. Adjoin $b$ of valuation $r$ and substitute $x=\exp(bY)$. Exponential and logarithm give inverse analytic bijections between $v(Y)\ge0$ and $B_r$, with nonvanishing derivatives. The power series of $P(\exp(bY))$ has coefficients $M_k(P)b^k/k!$. After normalization their valuations tend to $+\infty$. The Weierstrass degree of the series on the closed unit disk is its last minimum index and counts its zeros with multiplicity. Legendre's formula gives the expression for $w_r$. If $m_i$ is an initial unit, compare the minimum at $k$ to $w_r(i)$. Since $k>i$, the term $(r-r_0)(i-k)$ is strictly negative, which proves \eqref{digit} even for ramified coefficients.
\end{proof}

\begin{proof}[Proof of Theorem \ref{transfer}]
Match the $t$ initial moments exactly. The two Vandermonde matrices are invertible over $\Q$, so this defines the stated isomorphism and its inverse. For a nonzero pair normalize by the same initial moment. The congruence of the two recurrences implies by induction that $m_k(E)-m_k(F)\in p^M\mathcal O_K$ for every $k$. If a high index minimizes on either side, \eqref{digit} gives valuation strictly below $M$, so it has exactly the same valuation on the other side. The initial moments agree without a congruence. Applying this comparison in both directions proves equality of minima and of every minimizing index. For $M=0$ no high index can minimize and the same reasoning applies. The lemma gives the result at rational depths. Each polynomial has finitely many root depths and its finite depths are rational, so these counts determine all the others and the open-disk count.
\end{proof}

\section{An exact finite algorithm}
Choose the least nonnegative representatives $c$ modulo $q$. For each multiplicity $h_c$ replace the exponents in that residue by $c,c+q,\ldots,c+(h_c-1)q$. The resulting support $F$ lies in $[0,tq-1]$. Thus an arbitrary support, even with exponents of unbounded size, can be replaced by a uniformly bounded list. For all supports of size $t$, enumerate the $\binom{q+t-1}{t}$ residue multisets and take the largest answer. A smaller support can be padded with unused exponents.

Here is a simple elimination procedure for each bounded support. Work with the coefficient vectors of $(1+Y)^n$, $n\in F$. For a nonzero vector $a=(a_0,\ldots,a_L)$ and a closed depth $r$, let $\lambda(a)$ be the largest minimizing index for $v(a_j)+rj$. For the open disk use $r=r_0$ and the \emph{smallest} minimizing index. These two tie rules follow from the Newton polygon: on factoring into linear factors the least minimizing index counts roots strictly inside the radius, while the largest includes the boundary. More explicitly, the initial terms for a product multiply in the graded residue ring; the extreme minimizing indices therefore add, and the assertion for a linear polynomial is immediate. For irrational $r$ either take limits between adjacent finite breakpoints or apply the same factorwise rule.

Choose one column, record $\lambda$, then clear its coefficient at that index in every remaining column by ordinary rational column operations. Repeat until all $t$ columns have been selected. Previously cleared coefficients remain zero, hence the recorded indices are distinct. They give \emph{exactly} the $t$ possible counts, even after extending the coefficient field. Indeed, in any nontrivial linear combination of the resulting basis take the summands of lowest weighted coefficient valuation after scalar multiplication. Their distinguished indices are distinct. At the largest of these indices (smallest for the open rule) the distinguished coefficient cannot be canceled by another summand. All more extreme indices have strictly greater valuation. The extremal index of the sum is therefore one of the recorded values. Each pivot is itself attained by a rational basis vector. This proves, in particular, that there are precisely $t$ different attainable counts on each disk for a $t$-dimensional support space.

All operations use rational coefficients and, at rational depths, integer comparisons. The same calculation with $r$ symbolic gives a finite step function of the depth: for a bounded support, a fixed elimination path compares affine functions with integer intercepts and slopes, and there are finitely many pivot sequences. For a single depth, the canonical lift followed by elimination uses $O(t^2L)$ rational arithmetic operations for $L=\max F$. This operation count excludes bit complexity and the cost of enumerating all supports.

The original $E$ also has rational witnesses by moment matching. A rational polynomial yielding the desired count can be perturbed on its support to have all terms present and simple nonzero roots. For any fixed depth, the finitely many ordinary Taylor coefficients after $X=1+Y$ keep the same minimizing indices under sufficiently small coefficient perturbations: nonzero terms at the minimum have unchanged valuations, and all others remain above the minimum. This applies also to the open-disk tie rule. After factoring out $X^{\min E}$, fix a nonzero leading coefficient, clear denominators and add sufficiently $p$-divisible integers to the other allowed coefficients. The Chinese remainder theorem can simultaneously impose Eisenstein's conditions at a second prime away from the leading coefficient. This gives a separable polynomial over $\Q$. If a leading coefficient was initially zero, first perturb it to be nonzero. These steps concern a specified disk; Theorem~\ref{transfer} already matches all depths simultaneously with multiplicity.

\begin{center}
\begin{tabular}{ccll}
\toprule $p$ & $t$ & supports with the same residues modulo $p$ & open counts\\\midrule
$2$ & $8$ & $0,1,4,5,8,9,12,13$ & $0,1,\ldots,7$\\
    &     & $0,1,3,4,5,8,9,12$ & $0,1,\ldots,6,8$\\
$3$ & $9$ & $0,1,2,9,10,11,18,19,20$ & $0,1,\ldots,8$\\
    &     & $0,1,2,5,9,10,11,18,19$ & $0,1,\ldots,7,9$\\\bottomrule
\end{tabular}
\end{center}
In both pairs $M=2$. Exact elimination (see the supplied verifier) exhibits the need to retain more than one residue digit. This does not assert minimality of $M$ for arbitrary parameters.

\section{Sparse witnesses of prime-power degree}
\begin{lemma}\label{tower}
For every prime $p$ and integer $h\ge1$ there is an integer polynomial of degree $p^h$ with at most $p^{h-1}+2$ terms, all of whose roots are simple and in $B^\circ$. Its support lies in $\{0,p,2p,\ldots,p^h\}\cup\{1\}$.
\end{lemma}
\begin{proof}
Choose $\pi^{p-1}=p$, put $X=1+\pi Z$ and $H=(X^p-1)/\pi^p$. As a polynomial in $Z$, $H$ is integral, monic, and reduces to $Z+Z^p$ modulo the maximal ideal. Indeed, its linear coefficient is $1$, its leading coefficient is $1$, and its intermediate coefficients have positive valuation. Consider the following polynomial in $X$, using $Z=(X-1)/\pi$:
\[
 W_h=\pi^{p^h}\left(\sum_{j=0}^{h-1}(-1)^{h-1-j}H^{p^j}+(-1)^h Z\right).
\]
The powers $\pi^{p^h-p^{j+1}}$ and $\pi^{p^h-1}$ appearing after expansion are nonnegative integral powers of $p$. This proves the coefficient and support claims. It is monic. The normalized transform $\pi^{-p^h}W_h(1+\pi Z)$ reduces to $Z^{p^h}$ by telescoping in characteristic $p$. All its zeros are integral and reduce to zero, so the original roots have depth strictly greater than $r_0$. Adding an arbitrarily highly $p$-divisible integer constant preserves the reduction. Avoiding the finitely many critical values gives simple roots. At $h=2$ this specializes to the jump construction in \cite{repo}.
\end{proof}

\subsection{The first staircase after the Frobenius jump}

\begin{theorem}\label{staircase}
For every prime $p$ and every $p+2\le t\le2p+1$,
\[
 N_p^\circ(t)=p^2+\left\lfloor\frac{t-p-2}{2}\right\rfloor .
\]
Moreover $N_2^\circ(6)=8$.
\end{theorem}

\begin{proof}
Write $A=\lfloor(t-p-2)/2\rfloor$ and $S_p(k)$ for the base-$p$ digit sum.
Adding zero coefficients allows us to bound polynomials with exactly $t$ available exponents.
If all exponents occurring have the same residue, write the polynomial as
$X^u Q(X^{p^e})$, choosing $e$ maximal so at least two residues occur in $Q$. The map
$x\mapsto x^{p^e}$ is injective on the open disk by the $p$-adic logarithm, maps it into
itself, and has nonzero derivative there. We can now pad the support of $Q$ to $t$
distinct exponents. A monomial with one term contributes zero roots.

Take a rational closed depth $r=1/(p-1)+\delta$ with $\delta>0$, and put
$w(k)=\delta k+S_p(k)/(p-1)$. Normalize the initial moments as in Lemma~\ref{moment}, so that all $M_k$ are integral and some $M_i$ with $i<t$ is a unit.
For $i<t\le2p+1$ we have $S_p(i)\le p$, also when $p=2$.
If an index $k\ge t$ minimizes $w(k)+v_p(M_k)$, comparison with that $i$ gives
\[
 S_p(k)\le p-1,\qquad v_p(M_k)<1 .
\]
In $R=\mathbb F_p[T]/\prod_{n\in E}(T-\bar n)$ a relation expressing $T^k$ using
$T^j$ with $j<k$ and $S_p(j)\le S_p(k)$ excludes such a minimizer. Indeed its lift
gives $M_k$ as an integral linear combination of those $M_j$ modulo $p$,
so some $v_p(M_j)\le v_p(M_k)$ while $w(j)<w(k)$.

Suppose first that each residue has multiplicity at most $p$.
On the primary factor at $c$ all $T^{p^i}$, $i\ge1$, equal $c$.
If $a$ is the units digit of $k$ and $b$ the sum of its higher digits, then
$T^k=T^{a+pb}$ in $R$. Since $a+b\le p-1$, this replacement has equal digit sum
and index below $p^2$. It excludes every $k>p^2+A$.

Otherwise exactly one residue has multiplicity $m>p$. A monomial shift puts it at
zero, and the presence of another residue gives $m\le t-1\le2p$.
The product of the other primary factors, denoted by $R'$, has dimension $h=t-m$.
On a factor at $c\ne0$ in $R'$ we have $T^p=U=c$, an invertible element.
On $\mathbb F_p[T]/(T^m)$ every power of index at least $m$ vanishes.
Consider $k>p^2+A$. If $b\ge2$, the power $a+pb$ is at least $2p$ and below $p^2$,
so the replacement from the preceding paragraph still works on every factor.
If $b=1$, the element of $R'$ to be represented is $UT^a$.
When $a\ge m-p$ the power $T^{p+a}$ already represents it and vanishes on the zero
factor. Otherwise $a\le m-p-1$. Reducing $T^a$ modulo the monic degree-$h$
polynomial defining $R'$ expresses it in the powers $T^j$, $0\le j\le\min(a,h-1)$.
Each $UT^j$ is represented on the full algebra by $T^{p^2+j}$. Here
\[
 j\le\min(m-p-1,t-m-1)\le A,\qquad
 S_p(p^2+j)=1+j\le1+a=S_p(k).
\]
These lower indices exclude $k$ and give the upper bound at every closed depth,
hence on the open disk.

For sharpness, take the degree-$p^2$ polynomial $W$ furnished by
Lemma~\ref{tower}. Its $p^2$ roots are simple and lie in the open disk, and its
exponents lie in $\{0,p,2p,\ldots,p^2\}\cup\{1\}$. Choose $A$ distinct rational
integers $u_j\in1+p^2\mathbb Z$ outside the roots of $W$, with their $p$th
powers all distinct. Put
\[
 P(X)=W(X)\prod_{j=1}^{A}(X^p-u_j^p).
\]
Each added factor has exactly one root in the open disk, namely $u_j$.
Its other roots are $u_j\xi$ for nontrivial $p$th roots of unity $\xi$,
and $v_p(\xi-1)=1/(p-1)$ puts them on the boundary.
Thus all $p^2+A$ roots in the open disk are simple. The support of $P$
lies in the $p+A+1$ multiples $0,p,\ldots,(p+A)p$ and the $A+1$ integers
$1,1+p,\ldots,1+Ap$. Their total number is $p+2+2A\le t$.

For $p=2,t=6$, $\max_{i<6}S_2(i)=2$ again. A minimizing $k\ge6$ has digit sum one
and $v_2(M_k)<1$. After removing a common power as above all residual
multiplicities are at most five, so $T^{2^e}=T^8$ in the residue algebra for
$e\ge3$. This gives the upper bound eight. The degree-eight six-term example in
Lemma~\ref{tower} proves attainment.
\end{proof}

\section{Every fixed positive depth}
For completeness, finite determination extends to each region $r\ge r_*>0$, including depths outside the exponential chart. The following coarser precision uses Lenstra's effective uniform bound $B=B(p,t,r_*)$ for the number of geometric roots in $v(x-1)\ge r_*$ of any polynomial with at most $t$ terms \cite[Theorem 3.2 and Proposition 6.1]{len}. Enlarge $B$ to an integer at least $t-1$.

\begin{proposition}
Theorem~\ref{transfer} holds on every closed or open disk of depth $r\ge r_*$ if one uses $M_*=\max(1,v_p(B!))$ in place of $M$.
\end{proposition}
\begin{proof}
Match initial power moments and normalize as before. If $L(g)=\sum c_n g(n)$ denotes the normalized moment functional, the coefficient of $Y^j$ in $P(1+Y)$ is $L((T)_j)/j!$. Here $(T)_j$ is monic integral of degree $j$, so its initial values under $L$ for $j<t$ are related to the initial power moments by an integral unipotent transformation. At least one, say at $i<t$, is a unit and the matched initial values are equal. All $L((T)_j)$ are integral and congruent modulo $p^{M_*}$ for the matched supports. Every minimizing index at depth $r$ lies at most at $B$ by the uniform bound applied to the closed disk. If $k\ge t$ minimizes, comparison with $i$ gives
\[
 v(L((T)_k))\le r(i-k)+v(k!)-v(i!)<v(k!)\le v(B!).
\]
Thus again every possible minimum keeps its exact valuation on both sides. The two tie conventions prove the assertion. The same bounded-support reduction and ordinary Taylor elimination now apply.
\end{proof}
The small digit modulus of Theorem~\ref{transfer} uses $r>r_0$ essentially. At $r=r_0$, $X^p-1$ has all $p$ roots in the closed disk, whereas it has only one in $B^\circ$.

\paragraph{Scope and next questions.}
ProveIt's sparse-equations section obtains an exact finite-field rank test for $t\le p$ and the first exceptional value at $p+2$; it asks which additional residue data suffice beyond that range. Theorem~\ref{transfer} supplies an effective answer for all $t$ and the staircase evaluates an entire range in closed form. Earlier work of van der Poorten and Rumely bounds zeros of $p$-adic exponential polynomials \cite{vr}. Lenstra gives a degree-independent bound for geometric roots at every positive fixed depth \cite{len}. Avenda\~no and Krick study root maxima over a specified local field \cite{ak}. The distinguishing assertions here are the simultaneous profile isomorphism, the exact algorithm and the explicit post-jump values. The best modulus for all $(p,t)$ and a formula for later bands remain open directions. The present proofs use no hypothesis about $L$-functions.

\begin{thebibliography}{9}
\bibitem{repo} V. Reshetnikov, \emph{ProveIt, GeneralizedRiemannHypothesis}, article section \emph{Sparse equations for singular moduli}, in particular \texttt{smd:thm:rankrule}, \texttt{smd:thm:jump} and \texttt{smd:q:roots}. \href{https://github.com/VladimirReshetnikov/ProveIt/blob/41a284d6ef8eb16d783bfc685ab99cc6e1d28daa/Analysis/GeneralizedRiemannHypothesis/Article/sections/smd.tex}{Pinned source}.
\bibitem{vr} A. J. van der Poorten and R. S. Rumely, \emph{Zeros of $p$-adic exponential polynomials II}, J. London Math. Soc. (2) 36 (1987), 1--15. \href{https://doi.org/10.1112/jlms/s2-36.1.1}{Publisher's record}.
\bibitem{len} H. W. Lenstra, Jr., \emph{On the factorization of lacunary polynomials}, Number Theory in Progress, vol.~1, 1999, 277--292. \href{https://pub.math.leidenuniv.nl/~lenstrahw/PUBLICATIONS/1999b/art.pdf}{Author's copy}.
\bibitem{ak} M. Avenda\~no and T. Krick, \emph{Sharp bounds for the number of roots of univariate fewnomials}, J. Number Theory 131 (2011). \href{https://arxiv.org/abs/1008.4808}{arXiv:1008.4808}.
\end{thebibliography}
\end{document}
